\section{Chern--Simons 作用量与 $k$ 的整性}\label{sec:cs}

\subsection{这一章要解决什么}
本文只用了一个作用量。本章要回答四个问题：为什么不用 Yang--Mills？Chern--Simons 作用量
(1.3) 的三个性质从哪来？为什么耦合常数 $k$ 必须是整数？Wilson 线与配分函数 (1.5)(1.6)
怎么读？其中第三个问题原文只给了“必要且充分“一句 \src{355}，本章把必要性与充分性分开证。

\subsection{原文说了什么}
设 $M$ 是定向三维流形，$G$ 是紧 Lie 群，$A=A_i\dd x^i$ 是 $\mathfrak g$-值 1-形式。
规范变换与场强：
\begin{equation}
A_i\;\longrightarrow\;A_i-D_i\varepsilon,\qquad
F_{ij}=[D_i,D_j]=\partial_iA_j-\partial_jA_i+[A_i,A_j]. \tag{1.1}
\end{equation}
标准 Yang--Mills 作用量需要度规：
\begin{equation}
L_0=\int_M\sqrt g\,g^{ik}g^{jl}\Tr(F_{ij}F_{kl}). \tag{1.2}
\end{equation}
Witten 用的是
\begin{equation}
\label{eq:cs}
L=\frac{k}{4\pi}\int_M\Tr\Big(A\wedge \dd A+\frac23A\wedge A\wedge A\Big)
 =\frac{k}{8\pi}\int\dd^3x\,\varepsilon^{ijk}\Tr\Big(A_i\partial_jA_k+\frac23A_i[A_j,A_k]\Big).
 \tag{1.3}
\end{equation}
在大规范变换（绕数 $m$）下 $L\to L+\text{const}\cdot m$ \src{354}；要求 $\exp(\ii L)$ 单值，
故常数是 $2\pi$ 的整数倍，于是 $k\in\mathbb Z$ \src{355}。Wilson 线与配分函数：
\begin{equation}
W_R(C)=\Tr_R\,P\exp\int_CA_i\dd x^i, \tag{1.5}
\end{equation}
\begin{equation}
Z(M;C_i,R_i)=\int\cD A\,\exp(\ii L)\prod_iW_{R_i}(C_i). \tag{1.6}
\end{equation}
取 $G=SU(N)$、所有 $R_i$ 为定义表示，$S^3$ 上的 (1.6) 给出 Jones 多项式及其两变量推广，
两个变量是 $N$ 与 $k$（解析延拓为复值）；$G=SO(N)$ 给 Kauffman 多项式 \src{356}。

\subsection{原文没说的细节}
\begin{enumerate}[label=G2.\arabic*,leftmargin=*]
  \item 为什么 (1.2) 的两个因子性质不同：$\sqrt g\,\dd^3x$ 是体积元，换度规只改标度；
        $g^{ik}g^{jl}$ 把 $F$ 的指标缩并，换度规会改变不同分量之间的混合方式。
        前者可靠“不积分体积“绕开，后者绕不开。\ours
  \item $\delta L=\frac{k}{2\pi}\int\Tr(\delta A\wedge F)$ 的中间步骤原文跳过 \src{354}。
  \item 大规范变换下“平移 $2\pi$ 的整数倍“的常数到底是多少，原文写成 $\text{const}\cdot m$，
        没有写出 $2\pi k$。\ours
  \item “必要且充分“只有一句；充分性需要单独论证。
  \item $\Tr$ 的归一化：原文说“长根平方为 2 的归一化“使 $k\in\mathbb Z$ \src{355}，
        但 $SU(N)$ 定义表示的迹与这个归一化差一个整数因子，用哪个都自洽，但不能混用。
  \item 弱耦合是大 $k$ 还是小 $k$？原文在 \S2 才说清，本章先给出判断依据。
\end{enumerate}

\subsection{补出的推导}

\begin{derivation}[D-2.1：变分]
从 (1.3)。记 $A^3=A\wedge A\wedge A$。变分：
\[
\delta L=\frac{k}{4\pi}\int\Tr\Big(\delta A\wedge \dd A+A\wedge \dd\,\delta A
+\frac23(\delta A\wedge A\wedge A+A\wedge\delta A\wedge A+A\wedge A\wedge\delta A)\Big).
\]
第二项分部积分（$M$ 闭）：
$\int\Tr(A\wedge\dd\,\delta A)=-\int\Tr(\dd A\wedge\delta A)=\int\Tr(\delta A\wedge\dd A)$，
因为 $\Tr$ 循环且 $\delta A$ 是 1-形式（交换两个 1-形式要变号两次）。
后三项在迹下循环重排后给出：
\[
\frac23\int\Tr\big(\delta A\wedge(A\wedge A+A\wedge A+A\wedge A)\big)
=2\int\Tr(\delta A\wedge A\wedge A).
\]
合并：
\[
\delta L=\frac{k}{4\pi}\cdot 2\int\Tr(\delta A\wedge(\dd A+A\wedge A))
=\frac{k}{2\pi}\int\Tr(\delta A\wedge F).
\]
因此驻点方程是 $F=0$：平坦联络。这正是第 \ref{sec:phase} 讲相空间是“平坦联络模空间“的来源。
\end{derivation}

\begin{derivation}[D-2.2：大规范变换与绕数]
取 $g:M\to G$，$A\to A^g=g^{-1}Ag+g^{-1}\dd g$。直接计算给出
\begin{equation}
L(A^g)-L(A)=\frac{k}{4\pi}\int_M\dd\,\Tr(\dd g\,g^{-1}\wedge A)
+\frac{k}{12\pi}\int_M\Tr(g^{-1}\dd g)^3. \tag{D-2.1}
\end{equation}
第一项是恰当形式，在闭流形上积分为零。第二项只依赖 $g$ 的同伦类，且
\begin{equation}
m(g):=\frac{1}{24\pi^2}\int_M\Tr(g^{-1}\dd g)^3\in\mathbb Z \tag{D-2.2}
\end{equation}
是 $g$ 的绕数（映射度）。于是
\begin{equation}
L(A^g)-L(A)=\frac{k}{12\pi}\cdot 24\pi^2m=2\pi k\,m. \tag{D-2.3}
\end{equation}
这就是原文“平移 $2\pi$ 的整数倍“的确切形式。\ours
\end{derivation}

\begin{derivation}[D-2.3：必要性]
若 $k\in\mathbb Z$，则 $\exp(\ii(L+2\pi km))=\exp(\ii L)$，路径积分对同类大规范变换不变。
反过来，若 $k\notin\mathbb Z$，取绕数 $m=1$ 的 $g$，$\exp(\ii L)$ 变成 $\exp(\ii L)e^{2\pi\ii k}\ne\exp(\ii L)$，
配分函数没有定义。所以 $k\in\mathbb Z$ 是\textbf{必要}的。
\end{derivation}

\begin{derivation}[D-2.4：充分性（原文只说了一句）]
充分性有两层。
\begin{enumerate}[label=(\alph*),leftmargin=*]
  \item \textbf{同一平凡化下}：若 $k\in\mathbb Z$，则对任何 $g$，(D-2.3) 保持 $\exp(\ii L)$ 不变，
        所以 $\exp(\ii L)$ 是规范不变的函数。
  \item \textbf{不同平凡化之间}：$L$ 依赖 $A$ 在 $M$ 上的平凡化；换平凡化相当于一次大规范变换，
        由 (a) 只差 $2\pi k\cdot m$，$\exp(\ii L)$ 不变。故 $\exp(\ii L)$ 是 $M$ 上良定义的。
\end{enumerate}
两条合起来才是“充分“。原文的措辞是 “this condition is necessary and sufficient” \src{355}，
但只演示了必要性；充分性属于我们补的论证。\ours
\end{derivation}

\begin{derivation}[D-2.5：Wilson 线与 (1.6)]
$W_R(C)$ 的定义 (1.5) 只用联络沿曲线的和乐，不需要度规（曲线本身有参数化，但和乐与参数化无关）。
把它插进路径积分 (1.6)：$\cD A$ 是规范不变的测度（这一点原文未证，见第 \ref{sec:phase} 讲），
$W_R$ 规范不变，故 (1.6) 是链环 $(C_i,R_i)$ 的不变量。
$S^3$ 上取 $G=SU(N)$、全部 $R_i$ 为定义表示，算出的就是 Jones 多项式：
两个变量是 $N$ 与 $k$，经 $q=e^{2\pi\ii/(N+k)}$ 合成为一个 $q$（见第 \ref{sec:verlinde} 讲 (4.21)）。
\end{derivation}

\begin{keypoint}[弱耦合是大 $k$]
(1.3) 里 $k$ 在指数上：$\exp(\ii L)=\exp\big(\ii k\cdot\frac{1}{4\pi}\int\Tr(A\dd A+\tfrac23A^3)\big)$。
所以大 $k$ 对应快速振荡，可以用平稳相位展开；$1/k$ 才是耦合常数。
这与 Yang--Mills 的 $1/g^2$ 展开\textbf{方向相反}，不能把 $k$ 当“荷“。
原文 \S2 的弱耦合展开正是 $k\to\infty$。
\end{keypoint}

\subsection{与其他来源的对照}
\begin{center}
\begin{tabular}{@{}llll@{}}
\toprule
来源 & 作用量写法 & 与本文 $L$ 的关系 & $k$ 的整性 \\
\midrule
本文 \src{354} & $\dfrac{k}{4\pi}\int\Tr(A\dd A+\tfrac23A^3)$ & —— & $k\in\mathbb Z$ \\
河野 p.74 & $\mathrm{CS}=\dfrac{1}{8\pi^2}\int\Tr(A\dd A+\tfrac23A^3)$ & $L=2\pi k\,\mathrm{CS}$ & $k\in\mathbb Z$ \\
Moore §2 & $\dfrac{\kappa}{4\pi}\int\Tr(\cdots)$ & $\kappa=k$ & 整数或半整数（依群） \\
李思 5.4 & 预量子化 $\hbar$ & $\hbar=1/(2\pi k)$ & 整性 $[\omega]\in H^2(\mathbb Z)$ \\
上野 §3 & level $\ell=k$ & 同 & $\ell\in\mathbb Z$ \\
\bottomrule
\end{tabular}
\end{center}
换算规则与推导见附录 A；比较 $k$ 之前必须先对齐归一化，否则会差 $2\pi$ 或 $N$。

\subsection{陷阱与约定}
\begin{itemize}
  \item \textbf{$\Tr$ 的表示}：$SU(N)$ 定义表示 vs 长根平方 $=2$ 的归一化差一个因子 $N$；
        文中的“$k\in\mathbb Z$”是对某一固定归一化说的。换归一化，$k$ 的整性条件跟着换。
  \item \textbf{定向}：$L$ 依赖定向，$M$ 反定向时 $L\to-L$。所有不变量都按给定定向定义。
  \item \textbf{边界}：$M$ 有边界时 (1.3) 要加边界项才能有良好变分问题；本文只在闭流形上用 (1.3)，
        带边情形在 \S4 的手术里通过胶合处理。
\end{itemize}

\subsection{练习}
\begin{enumerate}[leftmargin=*]
  \item 逐步验证 D-2.1（提示：先证 $\int\Tr(A\wedge\dd\,\delta A)=\int\Tr(\delta A\wedge\dd A)$，
        注意 $\Tr$ 循环与 1-形式的交换符号）。
  \item 对 $G=SU(2)$，取 $g$ 为 $S^3\to SU(2)$ 的一次绕数映射，计算 (D-2.2) 并验证 $m=1$。
  \item 若 $k\in\mathbb Z+\tfrac12$，(D-2.3) 给 $\exp(\ii L)\to-\exp(\ii L)$。为什么这仍然可能给出
        一个“理论“？（提示：考虑 $M$ 是否可定向、以及是否允许自旋结构；本文不讨论这种情形。）
  \item 用 (1.6) 说明：为什么 $S^3$ 上取全部 $R_i$ 为定义表示给 Jones，而取 $SO(N)$ 给 Kauffman？
  \item 估算：$k\to\infty$ 时 (1.6) 的平稳相位近似给什么？（对照第 \ref{sec:conformal} 讲 $k\to\infty$ 的维数）
\end{enumerate}

\begin{boundary}
\begin{itemize}
  \item 原文证明：$k\in\mathbb Z$ 的必要性（绕数论证）；(1.3) 不含度规、变分给 $F=0$。
  \item 我们补：D-2.1 的完整变分、D-2.2 的系数 $2\pi k$、D-2.4 的充分性、D-2.5 的测度说明。
  \item 未证：$\cD A$ 与 (1.6) 的严格定义（本文是物理路径积分）；边界项的显式形式（留到第 \ref{sec:surgery} 讲）。
\end{itemize}
\end{boundary}
